\emph{Conventions.} $(\mathrm{div}\,f)_{s}$ is the sum of $f$ over the links leaving $s$ minus the sum over the links entering it, $\mathrm{grad}$ is its transpose (the value at the start of a link minus the value at its end), and $\mathrm{curl}$ is the signed sum around a plaquette, with $\mathrm{div}\,\mathrm{curl}^{T} = 0$. From $[E_{l}, W(f)] = f_{l}W(f)$: $[(\mathrm{div}\,E)_{s}, W(f)] = (\mathrm{div}\,f)_{s}W(f)$, and for any real combination $\Lambda = \sum_{l}\lambda_{l}E_{l}$, $[\Lambda, m\,W(f)] = (\sum_{l}\lambda_{l}f_{l})\,m\,W(f)$ for $m$ in the matter algebra, which commutes with the field's variables. In $\Cl(1,1)$: $P_{1}v = v$, $vP_{1} = 0$, so $[P_{1}, v] = v$, $[P_{1}, w] = -w$, $[P_{2}, v] = -v$, $[P_{2}, w] = w$, $[P_{1}, e_{2}] = [P_{1}, v - w] = v + w = e_{1}$ (\claimref{synthesis.set-and-clear-of-a-bit-generate-the-hyperbolic-plane}; the last is in the stored proof of Theorem \claimref{synthesis.field-slaved-to-an-operator-source-does-not-commute-at-spacelike-separation}).

\emph{(1).} $[G_{s}, e_{2}] = -g[\rho_{s}, e_{2}] = -g(\delta_{sa} - \delta_{sb})e_{1}$, since $\rho_{a} = P_{1}$ and $\rho_{b} = P_{2} = 1 - P_{1}$; so $[G_{s}, U] = -gq(\delta_{sa} - \delta_{sb})e_{1}$. For $\mathrm{div}\,f = \delta_{a} - \delta_{b}$: $[G_{s}, v\,W(gf)] = g(\mathrm{div}\,f)_{s}\,v\,W(gf) - g[\rho_{s}, v]\,W(gf) = g\bigl((\delta_{sa} - \delta_{sb}) - (\delta_{sa} - \delta_{sb})\bigr)v\,W(gf) = 0$, and likewise for $w\,W(-gf)$. $X_{f}^{2} = v^{2}W(2gf) + w^{2}W(-2gf) - vw - wv = -1$, so $(p + qX_{f})(p - qX_{f}) = p^{2} + q^{2} = 1$. $P_{1}$ commutes with every $G_{s}$, hence so does $O_{f}$, and $O_{f} = p^{2}P_{1} - q^{2}X_{f}P_{1}X_{f} + pq[P_{1}, X_{f}] = p^{2}P_{1} + q^{2}P_{2} + pq\,(v\,W(gf) + w\,W(-gf))$, using $X_{f}P_{1}X_{f} = -wW(-gf)\,P_{1}\,vW(gf) = -wv = -P_{2}$ and $[P_{1}, X_{f}] = vW(gf) + wW(-gf)$. Both $\chi$ and $\kappa$ have divergence $\delta_{a} - \delta_{b}$.

\emph{(2).} $[E_{l}, O_{f}] = pq\,\bigl(gf_{l}\,vW(gf) - gf_{l}\,wW(-gf)\bigr) = g\,pq\,f_{l}X_{f}$, and $X_{f} \neq 0$. With $f = \chi$ and $f = \kappa$ this gives the two cases. An element built from several Weyl operators, $\sum_{i}m_{i}W(gf_{i})$, has a commutator that is not a multiple of itself, and the statement is not made for it.

\emph{(3).} $\Pi = \mathrm{grad}\,(\mathrm{div}\,\mathrm{grad})^{-1}\mathrm{div}$ is the orthogonal projector onto gradients, symmetric, defined on functions whose divergence has zero total, which every divergence has. $E^{L}_{l} = \sum_{m}\Pi_{lm}E_{m}$, so $[E^{L}_{l}, m\,W(f)] = (\Pi f)_{l}\,m\,W(f)$ and $\Pi(g\chi) = g\kappa$: $[E^{L}_{l}, O_{\chi}] = g\,pq\,\kappa_{l}X_{\chi}$, and the transverse part takes the rest. On a state with $G_{s}\Psi = 0$ for all $s$, $E^{L}\Psi = g\,\mathrm{grad}\,(\mathrm{div}\,\mathrm{grad})^{-1}\rho\,\Psi = g\sum_{s}K(s)\rho_{s}\Psi$. The count of links with $\kappa_{l} \neq 0$ is computed in the stated box, with $\kappa$ obtained exactly from the graph Laplacian; on the six links of the bisector $\kappa$ vanishes by the reflection that exchanges $a$ and $b$.

\emph{(4).} A slaved field is $E_{l} = \epsilon_{l}$ with $\epsilon_{l}$ in the matter algebra and $\mathrm{div}\,\epsilon = g\rho$. With $A = 0$ at the start, two steps of the field equation of (7) give $E - h^{2}\,\mathrm{curl}^{T}\mathrm{curl}\,E$, which equals $E$ exactly when $\mathrm{curl}\,E = 0$ (pair $\mathrm{curl}^{T}\mathrm{curl}\,E$ with $E$), and then every later step leaves $E$ unchanged, $A$ staying a gradient. On the box a function on the links with no curl is a gradient (the plaquettes span the cycles of the grid), so the slaved field that the field equation leaves unchanged is $E = g\,\mathrm{grad}\,(\mathrm{div}\,\mathrm{grad})^{-1}\rho = g\sum_{s}K(s)\rho_{s}$. With $\rho_{a} = P_{1}$, $\rho_{b} = 1 - P_{1}$: $E_{l} = g\,(K_{l}(a) - K_{l}(b))\,P_{1} + \text{number} = g\,\kappa_{l}P_{1} + \text{number}$, so $[E_{l}, m] = g\,\kappa_{l}[P_{1}, m]$ for every matter element $m$, and for the bare detection $[P_{1}, O] = pq\,e_{2}$ by clause (1) of the stored model ($[J, O] = (c_{1} - c_{2})pq\,e_{2}$ with $J = P_{1}$). A matter element commutes with the charge exactly when $[P_{1}, m] = 0$. The source is built from the commuting path numbers as in Proposition \claimref{field-theory.operator-newtonian-coupling}.

\emph{(5).} $E_{l} = g\,\chi_{l}P_{1} + \text{number}_{l}$, the numbers with divergence $g\delta_{b}$ plus the background: its divergence is $g(\delta_{a} - \delta_{b})P_{1} + g\delta_{b}$ plus the background, which is $g\rho$; so it is a slaved field, and $[E_{l}, O] = g\,\chi_{l}\,pq\,e_{2}$. It has $\mathrm{curl}\,\chi \neq 0$, so by the argument of (4) it is not left unchanged; the count of links changed after two steps is computed. The field equation preserves the divergence, and the gradient part of a function on the links is determined by its divergence, so the gradient part of the evolved field is $g\kappa P_{1}$ plus a number at every step (computed at three steps as well).

\emph{(6).} A chain has no plaquette, so every function on its links is a gradient, $\Pi$ is the identity, and a function on the links is determined by its divergence up to a constant: the slaved field is unique up to a number. Then (4) gives $[E_{l}, m] = g\,\chi_{l}[P_{1}, m]$, zero outside the interval.

\emph{(7).} One step is the composition of two shears, $E \mapsto E + h\,\mathrm{curl}^{T}\mathrm{curl}\,A$ with $A$ fixed and $A \mapsto A - hE$ with $E$ fixed, each of which preserves the brackets of the pairs $(A_{l}, E_{l})$; it preserves $\mathrm{div}\,E$ because $\mathrm{div}\,\mathrm{curl}^{T} = 0$; and $(\mathrm{curl}^{T}\mathrm{curl}\,A)_{l}$ involves only links that share a plaquette with $l$, so $n$ steps forward or backward reach only links at distance at most $n$ in that relation. The quantum is at rest, so there is no current and the field equation is the same in every state of the matter. $E_{l}(n)$ is a real combination of the $E_{m}(n_{d})$ and $A_{m}(n_{d})$; the $A_{m}$ commute with $W(f)$, so $[E_{l}(n), W(f)(n_{d})] = c\,W(f)$ with $c$ the value at $(l, n)$ of the solution whose data at step $n_{d}$ are $E = f$, $A = 0$; with $f = \pm g\chi$ this gives $[E_{l}(n), O_{\chi}(n_{d})] = pq\,S_{l}(n)X_{\chi}$. $S_{l}(n) = 0$ beyond the cone by the locality of the step; the count is computed. For the slaved field: moving the quantum at step $0$ from the middle site $o$ to $a$ lays a unit of flux on the link from $o$ to $a$ on top of the Coulomb field of the quantum at $o$, which is static; likewise for $b$; the two strings differ by $\chi$, so $E^{(a)}(n) - E^{(b)}(n)$ is $g$ times the solution with data $\chi$, $0$ at step $0$, and each $E^{(i)}$ satisfies the constraint for the quantum at its site at every step. $E^{\mathrm{sl}} = E^{(b)} + (E^{(a)} - E^{(b)})P_{1}$ and the commutator follows as in (4). The count $878$ of $2750$ and the first step $4$ are computed (for $g = 1$; the count does not depend on $g \neq 0$). The solution with data at step $0$ on the string reaches at step $n$ only links at distance at most $n - 1$ from the string, so every event where it is nonzero lies in the forward cone of the source event; that each of the $878$ does, and lies outside the cone of the detection, is also computed. The operator in $E^{\mathrm{sl}}_{l}(n)$ is $P_{1}$, the charge at $a$, constant while the quantum is at rest; it is the charge laid down at the source event, and $[P_{1}, O] = pq\,e_{2}$ is the commutator of that charge with the detection of the same quantum at the same sites $n_{d}$ steps later. The mechanism is that of the stored model's clause (2): a retarded field is a function of the source operator at an earlier event, which does not commute with a later observable of the source.

\emph{Remark.} $S(n)$ solves the field equation and has the divergence of $g\chi$ at every step, so $N + SP_{1}$ is a slaved field obeying the field equation, and $[N_{l} + S_{l}(n)P_{1}, O] = S_{l}(n)\,pq\,e_{2}$; the vanishing of $S$ at the $2750$ events spacelike to the detection at step $5$ is the computation of (7), and the count $86$ of $2562$ for a detection at step $7$ is computed with the same $S$.

\emph{Reading.} The canonical pair on a link is not an element of a floor of the tower (\claimref{field-theory.no-canonical-commutation-relations-tower}). The statements about the gravitational gauge field are argued and use standard facts cited from memory, as the extraction note lists; no clause above depends on them.
